\section{Introduction}\label{sec:overview}

\begin{theorem}\label{thm:main}
Let $k$ be a field of characteristic zero. Let $L$ and
$M$ be finite-dimensional Lie algebras over $k$, and suppose that
$L$ is solvable. If there is an isomorphism of unital associative
$k$-algebras
\[
 \Phi:U_k(L)\longrightarrow U_k(M),
\]
then $L$ and $M$ are isomorphic as Lie algebras over $k$.
\end{theorem}

We identify a Lie algebra with its image in its enveloping algebra, using the
Poincar\'e--Birkhoff--Witt theorem. For a Lie ideal $J\subseteq g$,
the notation $U(g)J$ means the two-sided ideal it generates. Thus
$U(g)/U(g)J=U(g/J)$. Write $N(g)$
for the nilradical, and use the convention
\[
 \gamma_1N=N,\qquad \gamma_{r+1}N=[N,\gamma_rN].
\]

The proof separates the information supplied by the associative isomorphism
from the deformation argument needed to turn it into a Lie isomorphism.
First, we prove that $M$ is solvable and normalize $\Phi$ at the
augmentation. Locally nilpotent inner derivations then identify the ideal
$U(L)[N,N]$ intrinsically. In its metabelian quotient, locally
finite inner derivations identify both the polynomial algebra on
$N/[N,N]$ and the quotient $L/N$. Taking the linear part at the
augmentation gives an actual graded Lie isomorphism
\[
 \varphi_0:\gr_FL\longrightarrow\gr_FM,
 \qquad F^0L=L,\quad
 F^rL=\gamma_rN(L)\quad(r\geq1),
\]
whose enveloping map is the associated graded of the normalized $\Phi$.

Next, the corresponding Rees Lie algebras are two polynomial families with
the same graded special fibre. Their enveloping algebras are isomorphic by
a map equal to the identity at that fibre. After a separate homogenization
parameter $h$ is introduced, this map has a uniform Laurent bound on
each coefficient of the Rees parameter $t$. This precise bound allows its
use as a gauge transformation in the full Hochschild complex. Affine
formality reflects the resulting equivalence to the linear Poisson
structures of the two Lie families.

Finally, the first derivative at the origin of the Poisson equivalence
produces a Lie isomorphism over a Laurent-series field. Its graded part is
a regular arc specializing to $\varphi_0$. The image of the filtered
isomorphism scheme on the associated graded is a closed coset. Therefore
the fibre with mark $\varphi_0$ is nonempty. Its automorphism group is
unipotent, and an explicit finite Galois averaging argument descends a
marked isomorphism to the original coefficient field.

Only finitely many coefficients of the Lie brackets and of $\Phi$ and
$\Phi^{-1}$ are used to define the data. We can consequently carry out the
complex formality argument after embedding a finitely generated
characteristic-zero coefficient field into $\CC$, and descend the result
before extending back to $k$.

The order of the two formal operations is essential: one first completes in
$h$, then inverts $h$, and only afterwards completes in $t$.
The relevant completed coefficient spaces have the form
$X[[h]][h^{-1}][[t]]$. No lower bound in $h$ uniform in the
outer $t$-degree is required.
